% ============================================================================
% P-Center Problem – Document
% ============================================================================
\documentclass[11pt,a4paper]{article}

% ---- Packages ----
\usepackage[utf8]{inputenc}
\usepackage[T5]{fontenc}           % Vietnamese font encoding
\usepackage[vietnamese,english]{babel}
\usepackage{amsmath,amssymb,amsthm}
\usepackage{mathtools}
\usepackage{algorithm}
\usepackage{algpseudocode}
\usepackage{booktabs}
\usepackage{graphicx}
\usepackage{tikz}
\usetikzlibrary{positioning, arrows.meta, calc}
\usepackage{hyperref}
\usepackage{cleveref}
\usepackage[margin=2.5cm]{geometry}
\usepackage{enumitem}
\usepackage{caption}

% ---- Theorem environments ----
\newtheorem{definition}{Định nghĩa}
\newtheorem{example}{Ví dụ}

% ---- Custom commands ----
\DeclareMathOperator*{\argmin}{arg\,min}

% ============================================================================
\title{\textbf{Bài toán P-Center}\\[6pt]
\large Mô tả, Mô hình ILP, Mô hình SAT và Dữ liệu thực nghiệm}
\author{}
\date{}

\begin{document}
\maketitle
\tableofcontents
\newpage

% ====================================================================
\section{Mô tả bài toán}
\label{sec:description}
% ====================================================================

\subsection{Bối cảnh và ứng dụng}

Bài toán \emph{P-Center} (hay \emph{Vertex P-Center}) là một bài toán tối ưu tổ hợp kinh điển trong lĩnh vực \textbf{định vị cơ sở} (\textit{facility location}).
Bài toán được ứng dụng rộng rãi trong thực tế:
\begin{itemize}[nosep]
    \item \textbf{Dịch vụ khẩn cấp}: Đặt trạm cứu hỏa / cứu thương sao cho thời gian phản hồi tối đa là nhỏ nhất.
    \item \textbf{Viễn thông}: Đặt các trạm phát / server sao cho khoảng cách (hoặc độ trễ) lớn nhất tới người dùng được tối thiểu hóa.
    \item \textbf{Y tế}: Đặt các bệnh viện hoặc trung tâm tiêm chủng phục vụ một vùng dân cư.
    \item \textbf{Logistics}: Đặt kho hàng để giảm thiểu khoảng cách vận chuyển xa nhất.
\end{itemize}

\subsection{Định nghĩa hình thức}

\begin{definition}[Vertex P-Center Problem]
Cho đồ thị vô hướng có trọng số $G = (V, E, w)$ với $|V| = n$ đỉnh, trong đó $w: E \to \mathbb{R}_{\ge 0}$ là hàm trọng số cạnh.
Gọi $d(i,j)$ là khoảng cách đường đi ngắn nhất giữa hai đỉnh $i$ và $j$.
Cho số nguyên dương $p$ ($1 \le p \le n$).

Bài toán \textbf{Vertex P-Center} yêu cầu tìm tập $S \subseteq V$ với $|S| = p$ (gọi là các \textbf{center}) sao cho:
\[
    \min_{S \subseteq V,\; |S|=p}\; \max_{i \in V}\; \min_{j \in S}\; d(i,j)
\]

Nghĩa là, ta cần chọn $p$ đỉnh làm center để \textbf{tối thiểu hóa khoảng cách lớn nhất} từ bất kỳ đỉnh nào đến center gần nhất.
\end{definition}

Giá trị tối ưu $r^* = \max_{i \in V} \min_{j \in S^*} d(i,j)$ được gọi là \textbf{bán kính phủ} (\textit{covering radius}) hay \textbf{p-radius}.

\subsection{Độ phức tạp}

Bài toán P-Center là \textbf{NP-hard} với $p \ge 2$ trên đồ thị tổng quát~\cite{kariv1979algorithmic}.
Ngay cả phiên bản quyết định -- ``\emph{Có tồn tại tập $S$ với $|S| \le p$ sao cho bán kính phủ $\le R$?}'' -- cũng là \textbf{NP-complete}.
Vì vậy, các phương pháp giải chính xác thường dựa trên ILP, SAT/MaxSAT, hoặc branch-and-bound.


% ====================================================================
\section{Ví dụ minh họa}
\label{sec:example}
% ====================================================================

\begin{example}
Xét đồ thị $G$ gồm 6 đỉnh $V = \{1, 2, 3, 4, 5, 6\}$ với ma trận khoảng cách đường đi ngắn nhất cho trong \Cref{tab:dist_matrix}. Số center cần chọn: $p = 2$.
\end{example}

\begin{table}[h]
\centering
\caption{Ma trận khoảng cách $d(i,j)$ giữa các đỉnh.}
\label{tab:dist_matrix}
\begin{tabular}{c|cccccc}
\toprule
$d$ & 1 & 2 & 3 & 4 & 5 & 6 \\
\midrule
1 & 0 & 2 & 5 & 7 & 8 & 9 \\
2 & 2 & 0 & 3 & 5 & 6 & 7 \\
3 & 5 & 3 & 0 & 2 & 3 & 4 \\
4 & 7 & 5 & 2 & 0 & 1 & 2 \\
5 & 8 & 6 & 3 & 1 & 0 & 1 \\
6 & 9 & 7 & 4 & 2 & 1 & 0 \\
\bottomrule
\end{tabular}
\end{table}

\paragraph{Phân tích.}
\begin{itemize}[nosep]
    \item Nếu chọn $S = \{2, 5\}$:
    \begin{itemize}[nosep]
        \item Đỉnh 1: $\min(d(1,2), d(1,5)) = \min(2, 8) = 2$
        \item Đỉnh 2: $\min(d(2,2), d(2,5)) = \min(0, 6) = 0$
        \item Đỉnh 3: $\min(d(3,2), d(3,5)) = \min(3, 3) = 3$
        \item Đỉnh 4: $\min(d(4,2), d(4,5)) = \min(5, 1) = 1$
        \item Đỉnh 5: $\min(d(5,2), d(5,5)) = \min(6, 0) = 0$
        \item Đỉnh 6: $\min(d(6,2), d(6,5)) = \min(7, 1) = 1$
    \end{itemize}
    Bán kính phủ: $\max(2, 0, 3, 1, 0, 1) = 3$.

    \item Nếu chọn $S = \{1, 6\}$:
    \begin{itemize}[nosep]
        \item Khoảng cách lớn nhất: $\max(\min(0,9), \min(2,7), \min(5,4), \min(7,2), \min(8,1), \min(9,0)) = \max(0, 2, 4, 2, 1, 0) = 4$.
    \end{itemize}
    Bán kính phủ: $4 > 3$, không tốt bằng phương án trên.

    \item Nghiệm tối ưu: $S^* = \{2, 5\}$ với $r^* = 3$.
\end{itemize}

\begin{figure}[h]
\centering
\begin{tikzpicture}[
    node distance=1.8cm,
    vertex/.style={circle, draw, minimum size=8mm, font=\small},
    center/.style={circle, draw, fill=blue!25, minimum size=8mm, font=\small\bfseries},
    edge label/.style={font=\scriptsize, midway, auto}
]
    % Nodes
    \node[center] (2) {2};
    \node[vertex] (1) [left=2cm of 2] {1};
    \node[vertex] (3) [right=2cm of 2] {3};
    \node[vertex] (4) [right=2cm of 3] {4};
    \node[center] (5) [right=2cm of 4] {5};
    \node[vertex] (6) [right=2cm of 5] {6};

    % Edges with weights
    \draw (1) -- node[edge label, above] {2} (2);
    \draw (2) -- node[edge label, above] {3} (3);
    \draw (3) -- node[edge label, above] {2} (4);
    \draw (4) -- node[edge label, above] {1} (5);
    \draw (5) -- node[edge label, above] {1} (6);

    % Covering arcs
    \draw[dashed, blue!60, thick, -{Stealth}] (2) to[bend left=30] node[edge label, above, blue!80] {$d=2$} (1);
    \draw[dashed, blue!60, thick, -{Stealth}] (2) to[bend right=20] node[edge label, below, blue!80] {$d=3$} (3);
    \draw[dashed, red!60, thick, -{Stealth}] (5) to[bend left=20] node[edge label, above, red!80] {$d=1$} (4);
    \draw[dashed, red!60, thick, -{Stealth}] (5) to[bend right=30] node[edge label, below, red!80] {$d=1$} (6);
\end{tikzpicture}
\caption{Nghiệm $S^* = \{2, 5\}$ với bán kính phủ $r^* = 3$.
Các đỉnh center tô đậm; mũi tên nét đứt thể hiện phân bổ các đỉnh về center gần nhất.}
\label{fig:example}
\end{figure}


% ====================================================================
\section{Các ràng buộc chính}
\label{sec:constraints}
% ====================================================================

Bài toán P-Center bao gồm \textbf{bốn nhóm ràng buộc} cốt lõi:

\begin{enumerate}[label=\textbf{(C\arabic*)}]
    \item \textbf{Ràng buộc số lượng center (Cardinality constraint):}
    Phải chọn đúng $p$ đỉnh làm center.

    \item \textbf{Ràng buộc gán (Assignment constraint):}
    Mỗi đỉnh phải được gán cho đúng một center.

    \item \textbf{Ràng buộc liên kết (Linking constraint):}
    Một đỉnh chỉ có thể được gán cho một center nếu center đó đã được mở.

    \item \textbf{Ràng buộc khoảng cách (Distance / minimax constraint):}
    Khoảng cách từ mỗi đỉnh đến center được gán không vượt quá giá trị $z$ (bán kính phủ), đồng thời $z$ được tối thiểu hóa.
\end{enumerate}


% ====================================================================
\section{Mô hình ILP}
\label{sec:ilp}
% ====================================================================

\subsection{Tham số đầu vào}

\begin{itemize}[nosep]
    \item $V = \{1, 2, \dots, n\}$: tập đỉnh (vừa là điểm yêu cầu, vừa là vị trí tiềm năng đặt center).
    \item $d_{ij} \in \mathbb{R}_{\ge 0}$: khoảng cách đường đi ngắn nhất giữa đỉnh $i$ và đỉnh $j$.
    \item $p \in \mathbb{Z}_{>0}$: số center cần chọn.
\end{itemize}

\subsection{Biến quyết định}

\begin{itemize}[nosep]
    \item $y_j \in \{0, 1\}$, $\forall j \in V$: bằng 1 nếu đỉnh $j$ được chọn làm center, 0 ngược lại.
    \item $x_{ij} \in \{0, 1\}$, $\forall i, j \in V$: bằng 1 nếu đỉnh $i$ được gán cho center $j$, 0 ngược lại.
    \item $z \in \mathbb{R}_{\ge 0}$: bán kính phủ (khoảng cách lớn nhất).
\end{itemize}

\subsection{Mô hình}

\begin{align}
    \min \quad & z \label{eq:obj} \\[4pt]
    \text{s.t.} \quad
    & \sum_{j \in V} y_j = p \label{eq:cardinality} \\[4pt]
    & \sum_{j \in V} x_{ij} = 1 && \forall\, i \in V \label{eq:assign} \\[4pt]
    & x_{ij} \le y_j && \forall\, i, j \in V \label{eq:link} \\[4pt]
    & \sum_{j \in V} d_{ij}\, x_{ij} \le z && \forall\, i \in V \label{eq:dist} \\[4pt]
    & y_j \in \{0, 1\} && \forall\, j \in V \label{eq:binary_y} \\[4pt]
    & x_{ij} \in \{0, 1\} && \forall\, i, j \in V \label{eq:binary_x} \\[4pt]
    & z \ge 0 \label{eq:nonneg}
\end{align}

\subsection{Giải thích}

\begin{itemize}[nosep]
    \item \eqref{eq:obj}: Hàm mục tiêu -- tối thiểu hóa bán kính phủ $z$.
    \item \eqref{eq:cardinality}: \textbf{(C1)} -- Chọn đúng $p$ center.
    \item \eqref{eq:assign}: \textbf{(C2)} -- Mỗi đỉnh $i$ được gán cho đúng một center.
    \item \eqref{eq:link}: \textbf{(C3)} -- Đỉnh $i$ chỉ được gán cho $j$ nếu $j$ là center ($y_j = 1$).
    \item \eqref{eq:dist}: \textbf{(C4)} -- Khoảng cách từ đỉnh $i$ đến center được gán không vượt quá $z$.
    \item \eqref{eq:binary_y}--\eqref{eq:nonneg}: Miền giá trị biến.
\end{itemize}

\paragraph{Nhận xét.}
Mô hình có $O(n^2)$ biến nhị phân và $O(n^2)$ ràng buộc.
Trong thực tế, có thể cải thiện bằng cách:
\begin{itemize}[nosep]
    \item Loại bỏ các biến $x_{ij}$ khi $d_{ij} > z_{\text{UB}}$ (upper bound biết trước).
    \item Sử dụng các bất đẳng thức hợp lệ (\textit{valid inequalities}) và kỹ thuật \textit{branch-and-cut}.
\end{itemize}


% ====================================================================
\section{Mô hình SAT}
\label{sec:sat}
% ====================================================================

Thay vì giải bài toán tối ưu trực tiếp, cách tiếp cận SAT biến đổi P-Center thành \textbf{một chuỗi bài toán quyết định} thông qua tìm kiếm nhị phân trên bán kính $R$.

\subsection{Ý tưởng: Quy về bài toán Set Covering}

\paragraph{Bước 1: Sắp xếp các giá trị khoảng cách.}
Gọi $D = \{d_{ij} : i, j \in V, i \ne j\}$ là tập tất cả giá trị khoảng cách phân biệt, sắp xếp tăng dần: $d^{(1)} < d^{(2)} < \cdots < d^{(L)}$.

\paragraph{Bước 2: Tìm kiếm nhị phân.}
Sử dụng \textbf{tìm kiếm nhị phân} trên tập $D$ để tìm giá trị $R^*$ nhỏ nhất sao cho bài toán quyết định sau có lời giải:

\begin{quote}
\emph{``Có tồn tại tập $S \subseteq V$ với $|S| \le p$ sao cho mọi đỉnh $i \in V$ đều có ít nhất một center $j \in S$ với $d(i,j) \le R$?''}
\end{quote}

Bài toán quyết định này chính là bài toán \textbf{Set Covering}:
với mỗi đỉnh $j$, định nghĩa \emph{tập phủ}:
\[
    C_j(R) = \{i \in V : d(i,j) \le R\}
\]
Cần tìm $\le p$ tập $C_j$ sao cho hợp của chúng phủ toàn bộ $V$.

\paragraph{Bước 3: Mã hóa SAT.}
Bài toán Set Covering được mã hóa thành công thức CNF và giải bằng SAT solver.

\subsection{Mã hóa CNF}

\subsubsection{Biến Boolean}
\begin{itemize}[nosep]
    \item $y_j \in \{\texttt{true}, \texttt{false}\}$, $\forall j \in V$: $y_j = \texttt{true}$ nếu đỉnh $j$ được chọn làm center.
\end{itemize}

\subsubsection{Mệnh đề phủ (Coverage clauses)}

Với mỗi đỉnh $i \in V$, phải tồn tại ít nhất một center trong bán kính $R$:
\begin{equation}
    \bigvee_{j \in V:\; d(i,j) \le R} y_j \qquad \forall\, i \in V
    \label{eq:coverage}
\end{equation}

Mỗi ràng buộc \eqref{eq:coverage} tương ứng với \textbf{một mệnh đề} (clause) trong CNF.
Tổng cộng có $n$ mệnh đề phủ.

\subsubsection{Ràng buộc lực lượng: At-Most-$p$}

Cần đảm bảo $\sum_{j \in V} y_j \le p$, tức tối đa $p$ biến $y_j$ nhận giá trị \texttt{true}.

Ràng buộc này được mã hóa bằng \textbf{Sequential Counter} (Sinz, 2005~\cite{sinz2005towards}):

\paragraph{Biến phụ.}
Đặt $s_{j,k}$ ($j = 1, \dots, n$; $k = 1, \dots, p$) là biến phụ, trong đó $s_{j,k} = \texttt{true}$ nghĩa là ``trong số $y_1, \dots, y_j$ có ít nhất $k$ biến bằng \texttt{true}''.

\paragraph{Mệnh đề mã hóa.}
\begin{align}
    & \neg y_1 \lor s_{1,1} \label{eq:sc1} \\
    & \neg y_j \lor s_{j,1} && \forall\, j = 2, \dots, n \label{eq:sc2} \\
    & \neg s_{j-1,1} \lor s_{j,1} && \forall\, j = 2, \dots, n \label{eq:sc3} \\
    & \neg y_j \lor \neg s_{j-1,k-1} \lor s_{j,k} && \forall\, j = 2, \dots, n;\; k = 2, \dots, p \label{eq:sc4} \\
    & \neg s_{j-1,k} \lor s_{j,k} && \forall\, j = 2, \dots, n;\; k = 2, \dots, p \label{eq:sc5} \\
    & \neg y_j \lor \neg s_{j-1,p} && \forall\, j = 2, \dots, n \label{eq:sc6}
\end{align}

\begin{itemize}[nosep]
    \item \eqref{eq:sc1}--\eqref{eq:sc3}: Khởi tạo và truyền bit đầu tiên.
    \item \eqref{eq:sc4}--\eqref{eq:sc5}: Truyền đếm: nếu $y_j$ bằng true và đã có $k-1$ biến true, thì $s_{j,k}$ phải true.
    \item \eqref{eq:sc6}: \textbf{Mệnh đề cấm}: nếu đã có $p$ biến true ($s_{j-1,p}$), thì $y_j$ phải false $\Rightarrow$ đảm bảo at-most-$p$.
\end{itemize}

Mã hóa Sequential Counter thêm $O(np)$ biến phụ và $O(np)$ mệnh đề.

\subsection{Thuật toán tổng thể}

\begin{algorithm}[h]
\caption{Giải P-Center bằng SAT}
\label{alg:sat_pcenter}
\begin{algorithmic}[1]
\Require Đồ thị $G=(V,E,w)$, số center $p$
\Ensure Tập center tối ưu $S^*$, bán kính tối ưu $r^*$
\State Tính ma trận khoảng cách ngắn nhất $d_{ij}$, $\forall i,j \in V$
\State $D \gets$ tập giá trị $d_{ij}$ phân biệt, sắp xếp tăng dần
\State $lo \gets 1$, $hi \gets |D|$
\While{$lo < hi$}
    \State $mid \gets \lfloor (lo + hi) / 2 \rfloor$
    \State $R \gets D[mid]$
    \State Xây dựng CNF: mệnh đề phủ \eqref{eq:coverage} + at-most-$p$ \eqref{eq:sc1}--\eqref{eq:sc6}
    \If{SAT\_Solver(CNF) trả về SAT}
        \State $hi \gets mid$
        \State Lưu nghiệm $S$
    \Else
        \State $lo \gets mid + 1$
    \EndIf
\EndWhile
\State \Return $S^*$, $r^* = D[lo]$
\end{algorithmic}
\end{algorithm}

\paragraph{Phân tích.}
\begin{itemize}[nosep]
    \item Số lần gọi SAT solver: $O(\log L)$ với $L = |D| \le \binom{n}{2}$.
    \item Mỗi lần gọi, CNF có $n$ mệnh đề phủ $+$ $O(np)$ mệnh đề cardinality.
    \item Phương pháp này rất hiệu quả trong thực tế nhờ sức mạnh của các SAT solver hiện đại (CaDiCaL, Glucose, MiniSat, \ldots).
\end{itemize}


% ====================================================================
\section{Các phương pháp tối ưu hóa (Optimisation Approaches)}
\label{sec:optimisation}
% ====================================================================

Bài toán P-Center là bài toán \emph{tối ưu} (minimisation), nhưng SAT solver chỉ giải bài toán \emph{quyết định} (decision).
Do đó, cần một chiến lược bọc ngoài (\textit{outer loop}) để chuyển bài toán tối ưu thành chuỗi bài toán quyết định.

Phần này trình bày ba phương pháp chính: \textbf{Non-Incremental SAT}, \textbf{Incremental SAT}, và \textbf{MaxSAT}.


% -------------------------------------------------------------------
\subsection{Non-Incremental SAT (Giải SAT nhiều lần)}
\label{sec:non-incremental}
% -------------------------------------------------------------------

\subsubsection{Ý tưởng}

Đây là phương pháp đơn giản nhất: sử dụng \textbf{tìm kiếm nhị phân} trên tập giá trị bán kính $D = \{d^{(1)}, d^{(2)}, \dots, d^{(L)}\}$.
Tại mỗi bước, \textbf{xây dựng lại toàn bộ công thức CNF từ đầu} và gọi một SAT solver mới.

\subsubsection{Quy trình chi tiết}

\begin{enumerate}[nosep]
    \item \textbf{Tiền xử lý}: Tính ma trận $d_{ij}$, thu thập tập $D$ các giá trị khoảng cách phân biệt, sắp xếp tăng dần.
    \item \textbf{Tìm kiếm nhị phân}: Đặt $lo = 1$, $hi = |D|$.
    \item  Tại mỗi bước $mid = \lfloor(lo + hi)/2\rfloor$, đặt $R = D[mid]$:
    \begin{enumerate}[nosep, label=(\alph*)]
        \item Xây dựng \textbf{CNF mới hoàn toàn}:
            \begin{itemize}[nosep]
                \item $n$ mệnh đề phủ: $\bigvee_{j: d(i,j) \le R} y_j$ cho mỗi $i$.
                \item Mã hóa at-most-$p$ (ví dụ Sequential Counter): $O(np)$ mệnh đề.
            \end{itemize}
        \item Khởi tạo một \textbf{SAT solver mới} và nạp CNF.
        \item Gọi solver:
            \begin{itemize}[nosep]
                \item Nếu \textbf{SAT}: $hi \gets mid$, lưu nghiệm.
                \item Nếu \textbf{UNSAT}: $lo \gets mid + 1$.
            \end{itemize}
        \item \textbf{Hủy solver} (toàn bộ trạng thái nội bộ, learned clauses bị mất).
    \end{enumerate}
    \item Lặp lại cho đến $lo = hi$. Trả về $R^* = D[lo]$.
\end{enumerate}

\subsubsection{Pseudocode}

\begin{algorithm}[H]
\caption{Non-Incremental SAT cho P-Center}
\label{alg:non_incremental}
\begin{algorithmic}[1]
\Require Đồ thị $G=(V,E,w)$, số center $p$
\Ensure Tập center $S^*$, bán kính tối ưu $r^*$
\State Tính $d_{ij}$, $\forall i,j \in V$
\State $D \gets \text{sorted\_unique}(\{d_{ij}\})$
\State $lo \gets 1$, $hi \gets |D|$, $S^* \gets \emptyset$
\While{$lo < hi$}
    \State $mid \gets \lfloor (lo + hi) / 2 \rfloor$
    \State $R \gets D[mid]$
    \State $\varphi \gets \emptyset$ \Comment{Xây CNF mới từ đầu}
    \ForAll{$i \in V$}
        \State $\varphi \gets \varphi \cup \big\{  \bigvee_{j: d(i,j) \le R} y_j \big\}$ \Comment{Coverage clause}
    \EndFor
    \State $\varphi \gets \varphi \cup \textsc{SeqCounter}(y_1, \dots, y_n, p)$ \Comment{At-most-$p$}
    \State solver $\gets$ \textbf{new} SAT\_Solver()
    \State solver.add\_clauses($\varphi$)
    \If{solver.solve() $=$ SAT}
        \State $hi \gets mid$
        \State $S^* \gets \{j : \text{solver.value}(y_j) = \texttt{true}\}$
    \Else
        \State $lo \gets mid + 1$
    \EndIf
    \State \textbf{delete} solver \Comment{Hủy solver, mất toàn bộ learned clauses}
\EndWhile
\State \Return $S^*$, $r^* = D[lo]$
\end{algorithmic}
\end{algorithm}

\subsubsection{Phân tích}

\begin{itemize}[nosep]
    \item \textbf{Số lần gọi solver}: $O(\log L)$, với $L \le \binom{n}{2}$.
    \item \textbf{Chi phí mỗi lần}: Xây dựng CNF $O(n + np)$ + giải SAT (exponential worst case, nhưng thường nhanh trong thực tế).
    \item \textbf{Nhược điểm chính}: Mỗi lần gọi solver đều bắt đầu từ đầu -- \textbf{mất toàn bộ learned clauses} từ lần giải trước. Đây là lãng phí lớn vì các CNF giữa hai bước liên tiếp thường rất giống nhau (chỉ khác ở các mệnh đề phủ).
    \item \textbf{Ưu điểm}: Đơn giản, dễ cài đặt, không phụ thuộc API đặc thù của solver.
\end{itemize}


% -------------------------------------------------------------------
\subsection{Incremental SAT (SAT gia tăng)}
\label{sec:incremental}
% -------------------------------------------------------------------

\subsubsection{Ý tưởng}

Quan sát then chốt: trong quá trình tìm kiếm nhị phân, \textbf{ràng buộc cardinality (at-most-$p$) không thay đổi} -- chỉ có các \textbf{mệnh đề phủ thay đổi} khi $R$ thay đổi.
Incremental SAT \textbf{giữ nguyên solver} xuyên suốt quá trình, tận dụng:
\begin{itemize}[nosep]
    \item \textbf{Learned clauses}: Các mệnh đề học được từ lần giải trước vẫn hợp lệ (vì chúng là hệ quả logic của các mệnh đề cố định).
    \item \textbf{Heuristic state}: Thứ tự biến (VSIDS scores), phase saving, \ldots
\end{itemize}

\subsubsection{Cơ chế Assumption Literals}

SAT solver hiện đại hỗ trợ cơ chế \textbf{assumption} (giả thiết) qua giao diện IPASIR~\cite{balyo2016ipasir}:

\begin{itemize}[nosep]
    \item \texttt{solver.add\_clause(clause)}: Thêm mệnh đề \textbf{vĩnh viễn} (không thể gỡ bỏ).
    \item \texttt{solver.solve(assumptions=[a1, a2, ...])}: Giải dưới giả thiết các literal $a_1, a_2, \dots$ phải đúng. Khi lần gọi kết thúc, các assumption được \textbf{tự động gỡ bỏ}.
\end{itemize}

\paragraph{Kỹ thuật.}
Với mỗi giá trị bán kính $R_t$ (ứng với bước $t$ trong tìm kiếm nhị phân), tạo một \textbf{assumption literal} $\alpha_t$.
Các mệnh đề phủ tại bước $t$ được ``bảo vệ'' bởi $\alpha_t$:

\begin{equation}
    \neg \alpha_t \lor \bigvee_{j: d(i,j) \le R_t} y_j \qquad \forall\, i \in V
    \label{eq:guarded_coverage}
\end{equation}

Khi gọi solver với assumption $\alpha_t = \texttt{true}$, các mệnh đề này được kích hoạt.
Ở bước tiếp theo, $\alpha_t$ không được giả thiết nữa nên các mệnh đề cũ trở thành \textbf{tautology} (luôn đúng) và không ảnh hưởng đến solver.

\subsubsection{Cấu trúc CNF}

\begin{itemize}[nosep]
    \item \textbf{Mệnh đề cố định} (thêm một lần duy nhất):
    \begin{itemize}[nosep]
        \item At-most-$p$: Sequential Counter \eqref{eq:sc1}--\eqref{eq:sc6} trên $y_1, \dots, y_n$.
    \end{itemize}

    \item \textbf{Mệnh đề có điều kiện} (thêm tại mỗi bước $t$):
    \begin{itemize}[nosep]
        \item Coverage guarded: $(\neg \alpha_t \lor y_{j_1} \lor y_{j_2} \lor \cdots)$ cho mỗi $i$ với $\{j_1, j_2, \dots\} = \{j: d(i,j) \le R_t\}$.
    \end{itemize}

    \item \textbf{Gọi solver}: \texttt{solver.solve(assumptions=[$\alpha_t$])}.
\end{itemize}

\subsubsection{Pseudocode}

\begin{algorithm}[H]
\caption{Incremental SAT cho P-Center}
\label{alg:incremental}
\begin{algorithmic}[1]
\Require Đồ thị $G=(V,E,w)$, số center $p$
\Ensure Tập center $S^*$, bán kính tối ưu $r^*$
\State Tính $d_{ij}$, $\forall i,j \in V$
\State $D \gets \text{sorted\_unique}(\{d_{ij}\})$
\State solver $\gets$ \textbf{new} SAT\_Solver() \Comment{Khởi tạo \textbf{một lần duy nhất}}
\State solver.add\_clauses($\textsc{SeqCounter}(y_1, \dots, y_n, p)$) \Comment{Mệnh đề cố định}
\State $lo \gets 1$, $hi \gets |D|$, $S^* \gets \emptyset$
\While{$lo < hi$}
    \State $mid \gets \lfloor (lo + hi) / 2 \rfloor$
    \State $R \gets D[mid]$
    \State $\alpha \gets$ fresh\_variable() \Comment{Assumption literal mới}
    \ForAll{$i \in V$}
        \State solver.add\_clause$(\neg \alpha \lor \bigvee_{j: d(i,j) \le R} y_j)$ \Comment{Guarded coverage}
    \EndFor
    \If{solver.solve(assumptions=$[\alpha]$) $=$ SAT}
        \State $hi \gets mid$
        \State $S^* \gets \{j : \text{solver.value}(y_j) = \texttt{true}\}$
    \Else
        \State $lo \gets mid + 1$
    \EndIf
    \Comment{Solver giữ nguyên learned clauses \& heuristics}
\EndWhile
\State \Return $S^*$, $r^* = D[lo]$
\end{algorithmic}
\end{algorithm}

\subsubsection{Phân tích}

\begin{itemize}[nosep]
    \item \textbf{Solver chỉ khởi tạo một lần}: Tiết kiệm chi phí khởi tạo.
    \item \textbf{Learned clauses được tích lũy}: Solver ngày càng ``thông minh hơn'' qua các bước.
    \item \textbf{Mệnh đề cố định chiếm phần lớn CNF}: At-most-$p$ encoding ($O(np)$ mệnh đề) chỉ thêm một lần; mỗi bước chỉ thêm $n$ mệnh đề phủ mới.
    \item \textbf{Chi phí phụ}: Tổng cộng thêm $O(n \log L)$ mệnh đề guarded vào solver (không gỡ bỏ được, nhưng trở thành tautology).
    \item \textbf{Yêu cầu}: SAT solver hỗ trợ API incremental (IPASIR): CaDiCaL, Glucose, MiniSat, Lingeling, \ldots
\end{itemize}

\paragraph{Tại sao learned clauses vẫn hợp lệ?}
Khi giả thiết $\alpha_t$ bị gỡ, mệnh đề $(\neg \alpha_t \lor \ell_1 \lor \cdots \lor \ell_k)$ trở thành tautology (đúng với mọi assignment).
Bất kỳ mệnh đề nào suy ra (\textit{resolve}) từ tập mệnh đề gốc $+$ tautology đều vẫn là hệ quả logic hợp lệ.
Do đó, \textbf{tất cả learned clauses đều an toàn để giữ lại}.


% -------------------------------------------------------------------
\subsection{MaxSAT (Tối ưu hóa trực tiếp)}
\label{sec:maxsat}
% -------------------------------------------------------------------

\subsubsection{Ý tưởng}

Thay vì giải nhiều bài toán quyết định, MaxSAT cho phép \textbf{mã hóa trực tiếp mục tiêu tối ưu} vào một công thức duy nhất.
Ta sử dụng \textbf{Weighted Partial MaxSAT} (WPMS)~\cite{li2021maxsat}:
\begin{itemize}[nosep]
    \item \textbf{Hard clauses}: Phải thỏa mãn (vi phạm $\Rightarrow$ nghiệm không hợp lệ).
    \item \textbf{Soft clauses}: Có thể vi phạm, mỗi vi phạm chịu \textbf{trọng số} (weight). Mục tiêu: tối thiểu hóa tổng trọng số vi phạm.
\end{itemize}

\subsubsection{Mã hóa bài toán P-Center}

\paragraph{Biến Boolean.}
\begin{itemize}[nosep]
    \item $y_j$, $\forall j \in V$: đỉnh $j$ được chọn làm center.
    \item $t_k$, $k = 1, \dots, L$: biến ``ngưỡng'' (\textit{threshold}), trong đó $t_k = \texttt{true}$ nghĩa là bán kính tối ưu $\le d^{(k)}$.
\end{itemize}

\paragraph{Hard clauses} (phải thỏa mãn).

\begin{enumerate}[nosep, label=(H\arabic*)]
    \item \textbf{At-most-$p$}: Mã hóa $\sum_j y_j \le p$ bằng Sequential Counter (như §\ref{sec:sat}).
    \label{hc:cardinality}

    \item \textbf{Liên kết ngưỡng -- phủ}: Nếu $t_k = \texttt{true}$ (bán kính $\le d^{(k)}$), thì mọi đỉnh phải được phủ bởi ít nhất một center trong phạm vi $d^{(k)}$:
    \begin{equation}
        \neg t_k \lor \bigvee_{j: d(i,j) \le d^{(k)}} y_j \qquad \forall\, i \in V,\; k = 1, \dots, L
        \label{eq:maxsat_coverage}
    \end{equation}
    \label{hc:coverage}

    \item \textbf{Đơn điệu ngưỡng}: Nếu bán kính $\le d^{(k)}$ thì cũng $\le d^{(k+1)}$:
    \begin{equation}
        \neg t_k \lor t_{k+1} \qquad \forall\, k = 1, \dots, L-1
        \label{eq:monotone}
    \end{equation}
    \label{hc:monotone}
\end{enumerate}

\paragraph{Soft clauses} (tối ưu hóa).

Mỗi biến $t_k$ là soft clause đơn vị $\{t_k\}$ với trọng số $w_k$:
\begin{equation}
    \text{soft clause: } (t_k) \quad \text{weight } w_k \qquad \forall\, k = 1, \dots, L
    \label{eq:soft}
\end{equation}

Nếu $t_k$ vi phạm (tức $t_k = \texttt{false}$, nghĩa là bán kính $> d^{(k)}$), ta chịu chi phí $w_k$.

\paragraph{Chọn trọng số.}

Hai cách chọn trọng số phổ biến:
\begin{itemize}[nosep]
    \item \textbf{Trọng số đồng nhất}: $w_k = 1$, $\forall k$.
    MaxSAT solver sẽ tối đa số $t_k$ đúng, tức \textbf{tối thiểu $k^*$} nhỏ nhất mà $t_{k^*} = \texttt{false}$.
    Nhờ ràng buộc đơn điệu \eqref{eq:monotone}, nghiệm sẽ có dạng $t_1 = \cdots = t_{k^*-1} = \texttt{false}$, $t_{k^*} = \cdots = t_L = \texttt{true}$, tương ứng bán kính tối ưu $r^* = d^{(k^*)}$.

    \item \textbf{Trọng số lũy thừa}: $w_k = 2^{L-k}$.
    Ưu tiên mạnh việc thỏa mãn các $t_k$ nhỏ (bán kính nhỏ), đảm bảo solver tìm $k^*$ nhỏ nhất trước.
\end{itemize}

\subsubsection{Pseudocode}

\begin{algorithm}[H]
\caption{Weighted Partial MaxSAT cho P-Center}
\label{alg:maxsat}
\begin{algorithmic}[1]
\Require Đồ thị $G=(V,E,w)$, số center $p$
\Ensure Tập center $S^*$, bán kính tối ưu $r^*$
\State Tính $d_{ij}$, $\forall i,j \in V$
\State $D \gets \text{sorted\_unique}(\{d_{ij}\})$,\; $L \gets |D|$
\State solver $\gets$ \textbf{new} MaxSAT\_Solver()
\Statex
\State \Comment{\textbf{Hard clauses}}
\State solver.add\_hard($\textsc{SeqCounter}(y_1, \dots, y_n, p)$) \Comment{At-most-$p$}
\ForAll{$k = 1, \dots, L$}
    \ForAll{$i \in V$}
        \State solver.add\_hard$(\neg t_k \lor \bigvee_{j: d(i,j) \le D[k]} y_j)$ \Comment{Coverage}
    \EndFor
\EndFor
\ForAll{$k = 1, \dots, L-1$}
    \State solver.add\_hard$(\neg t_k \lor t_{k+1})$ \Comment{Monotonicity}
\EndFor
\Statex
\State \Comment{\textbf{Soft clauses}}
\ForAll{$k = 1, \dots, L$}
    \State solver.add\_soft$(t_k,\; w_k)$ \Comment{Tối đa hóa số $t_k$ đúng}
\EndFor
\Statex
\State solver.solve()
\State $S^* \gets \{j : \text{solver.value}(y_j) = \texttt{true}\}$
\State $k^* \gets \min\{k : \text{solver.value}(t_k) = \texttt{true}\}$
\State \Return $S^*$, $r^* = D[k^*]$
\end{algorithmic}
\end{algorithm}

\subsubsection{Phân tích}

\begin{itemize}[nosep]
    \item \textbf{Một lần gọi solver duy nhất}: Không cần outer loop (binary search).
    \item \textbf{Kích thước công thức}: $O(np)$ hard clauses (cardinality) $+$ $O(nL)$ hard clauses (coverage) $+$ $L$ soft clauses $+$ $L$ monotonicity clauses.
    \item \textbf{Ưu điểm}: MaxSAT solver (RC2, Open-WBO, EvalMaxSAT, \ldots) có các kỹ thuật tối ưu nội bộ mạnh mẽ (stratification, core-guided search, \ldots) giúp không cần duyệt qua tất cả $L$ giá trị.
    \item \textbf{Nhược điểm}: Kích thước CNF lớn hơn đáng kể ($O(nL)$ vs $O(n)$ cho mỗi bước SAT), có thể chậm hơn trên instance lớn nếu $L$ lớn.
\end{itemize}


% -------------------------------------------------------------------
\subsection{So sánh các phương pháp}
\label{sec:comparison}
% -------------------------------------------------------------------

\begin{table}[h]
\centering
\caption{So sánh ba phương pháp tối ưu hóa dựa trên SAT cho P-Center.}
\label{tab:comparison}
\small
\begin{tabular}{p{3cm} p{3.8cm} p{3.8cm} p{3.8cm}}
\toprule
\textbf{Tiêu chí} & \textbf{Non-Incremental} & \textbf{Incremental SAT} & \textbf{MaxSAT} \\
\midrule
\textbf{Số lần gọi solver} & $O(\log L)$ lần, mỗi lần solver mới & $O(\log L)$ lần, cùng solver & 1 lần duy nhất \\
\addlinespace
\textbf{Learned clauses} & Mất sau mỗi bước & Tích lũy xuyên suốt & Nội bộ MaxSAT solver \\
\addlinespace
\textbf{Kích thước CNF mỗi bước} & $O(n + np)$ & $O(n)$ mới $+$ $O(np)$ cố định & $O(nL + np)$ tổng \\
\addlinespace
\textbf{Yêu cầu API} & API cơ bản & IPASIR (assumptions) & WCNF (hard/soft) \\
\addlinespace
\textbf{Solver phù hợp} & MiniSat, CaDiCaL, Glucose & CaDiCaL, Lingeling, Glucose & RC2, Open-WBO, EvalMaxSAT \\
\addlinespace
\textbf{Độ phức tạp cài đặt} & Đơn giản & Trung bình & Trung bình \\
\addlinespace
\textbf{Hiệu quả thực tế} & Tốt cho instance nhỏ & Tốt nhất cho instance trung bình -- lớn & Tốt khi $L$ không quá lớn \\
\bottomrule
\end{tabular}
\end{table}

\paragraph{Hướng dẫn chọn phương pháp.}
\begin{itemize}[nosep]
    \item Nếu cần \textbf{cài đặt nhanh, đơn giản}: dùng \textbf{Non-Incremental}.
    \item Nếu muốn \textbf{hiệu quả tốt nhất} trên instance lớn: dùng \textbf{Incremental SAT} (tận dụng learned clauses).
    \item Nếu muốn \textbf{formulation gọn, một lần gọi}: dùng \textbf{MaxSAT} (nhưng cần cân nhắc kích thước $L$).
\end{itemize}


% ====================================================================
\section{Dữ liệu thực nghiệm (Benchmark Datasets)}
\label{sec:datasets}
% ====================================================================

Các nghiên cứu về P-Center thường đánh giá trên hai bộ dữ liệu chuẩn chính:

\subsection{Bộ pmed (OR-Library)}

\begin{itemize}[nosep]
    \item \textbf{Nguồn}: OR-Library, xây dựng bởi Beasley~\cite{beasley1990or}.
    \item \textbf{Nội dung}: 40 instance (\texttt{pmed1} -- \texttt{pmed40}).
    \item \textbf{Quy mô}: Số đỉnh $n$ từ 100 đến 900; giá trị $p$ từ 5 đến 200.
    \item \textbf{Cấu trúc}: Đồ thị thưa (sparse graph), khoảng cách giữa các cặp đỉnh được tính bằng đường đi ngắn nhất (thuật toán Floyd-Warshall hoặc Dijkstra).
    \item \textbf{Đặc điểm}:
    \begin{itemize}[nosep]
        \item Ban đầu thiết kế cho bài toán \textit{p-median}, nhưng được sử dụng phổ biến cho p-center vì cùng cấu trúc đồ thị.
        \item Nghiệm tối ưu đã được biết cho tất cả 40 instance.
    \end{itemize}
    \item \textbf{URL}: \url{https://people.brunel.ac.uk/~mastjjb/jeb/orlib/pmedinfo.html}
\end{itemize}

\subsection{Bộ TSPLIB}

\begin{itemize}[nosep]
    \item \textbf{Nguồn}: TSPLIB, xây dựng bởi Reinelt~\cite{reinelt1991tsplib}.
    \item \textbf{Nội dung}: Các instance gốc dùng cho bài toán TSP, được chuyển đổi cho p-center bằng cách tính khoảng cách Euclid giữa các điểm.
    \item \textbf{Quy mô}: Từ vài trăm đến hàng nghìn đỉnh (ví dụ: \texttt{att48}, \texttt{eil76}, \texttt{lin318}, \texttt{rat783}, \texttt{pr1002}, \ldots).
    \item \textbf{Đặc điểm}:
    \begin{itemize}[nosep]
        \item Đồ thị đầy đủ (complete graph) với khoảng cách Euclid 2D.
        \item Khó hơn bộ pmed do đồ thị đầy đủ $\Rightarrow$ ma trận khoảng cách dày đặc.
    \end{itemize}
    \item \textbf{URL}: \url{http://comopt.ifi.uni-heidelberg.de/software/TSPLIB95/}
\end{itemize}

\subsection{Bộ dữ liệu bổ sung}

Ngoài hai bộ chuẩn trên, một số nghiên cứu gần đây đã đề xuất các instance lớn hơn.
Ví dụ, Liu et al.~\cite{liu2020effective} thử nghiệm trên cả pmed và TSPLIB, sử dụng phương pháp quy về Set Covering $+$ SAT solver để giải hiệu quả các instance cỡ lớn.

\begin{table}[h]
\centering
\caption{Tổng hợp bộ dữ liệu benchmark cho P-Center.}
\label{tab:datasets}
\begin{tabular}{llccc}
\toprule
\textbf{Bộ dữ liệu} & \textbf{Loại đồ thị} & \textbf{Số instance} & \textbf{Quy mô $n$} & \textbf{Tham khảo} \\
\midrule
pmed (OR-Library) & Thưa (sparse) & 40 & 100 -- 900 & \cite{beasley1990or} \\
TSPLIB & Đầy đủ (Euclid 2D) & nhiều & 48 -- 3038+ & \cite{reinelt1991tsplib} \\
\bottomrule
\end{tabular}
\end{table}


% ====================================================================
% References
% ====================================================================
\bibliographystyle{plain}
\begin{thebibliography}{10}

\bibitem{kariv1979algorithmic}
O.~Kariv and S.~L.~Hakimi.
\newblock An algorithmic approach to network location problems. {I}: The $p$-centers.
\newblock \emph{SIAM Journal on Applied Mathematics}, 37(3):513--538, 1979.

\bibitem{beasley1990or}
J.~E.~Beasley.
\newblock {OR-Library}: distributing test problems by electronic mail.
\newblock \emph{Journal of the Operational Research Society}, 41(11):1069--1072, 1990.

\bibitem{reinelt1991tsplib}
G.~Reinelt.
\newblock {TSPLIB} --- a traveling salesman problem library.
\newblock \emph{ORSA Journal on Computing}, 3(4):376--384, 1991.

\bibitem{sinz2005towards}
C.~Sinz.
\newblock Towards an optimal {CNF} encoding of boolean cardinality constraints.
\newblock In \emph{Proceedings of the 11th International Conference on Principles and Practice of Constraint Programming (CP 2005)}, pages 827--831, 2005.

\bibitem{liu2020effective}
X.~Liu, Y.~Fang, J.~Chen, and Z.~L{\"u}.
\newblock Effective approaches to solve {P-Center} problem via set covering and {SAT}.
\newblock \emph{arXiv preprint arXiv:2008.xxxxx}, 2020.

\bibitem{minieka1970m}
E.~Minieka.
\newblock The $m$-center problem.
\newblock \emph{SIAM Review}, 12(1):138--139, 1970.

\bibitem{daskin1995network}
M.~S.~Daskin.
\newblock \emph{Network and Discrete Location: Models, Algorithms, and Applications}.
\newblock John Wiley \& Sons, 1995.

\bibitem{calik2015pcenter}
H.~\c{C}al{\i}k, M.~Labb\'e, and H.~Yaman.
\newblock p-center problems.
\newblock In G.~Laporte, S.~Nickel, and F.~Saldanha da Gama, editors, \emph{Location Science}, chapter~4, pages 79--92. Springer, 2015.

\bibitem{balyo2016ipasir}
T.~Balyo, A.~Biere, M.~Iser, and C.~Sinz.
\newblock {SAT Race 2015}.
\newblock \emph{Artificial Intelligence}, 241:45--65, 2016.
\newblock (Defines the IPASIR incremental SAT interface.)

\bibitem{li2021maxsat}
C.~M.~Li and F.~Manya.
\newblock {MaxSAT}, hard and soft constraints.
\newblock In A.~Biere, M.~Heule, H.~van~Maaren, and T.~Walsh, editors, \emph{Handbook of Satisfiability}, chapter~19, pages 903--927. IOS Press, 2nd edition, 2021.

\bibitem{een2003temporal}
N.~E{\'e}n and N.~S{\"o}rensson.
\newblock An extensible {SAT}-solver.
\newblock In \emph{SAT 2003}, LNCS 2919, pages 502--518. Springer, 2003.
\newblock (MiniSat -- foundational incremental SAT solver.)

\end{thebibliography}

\end{document}
